% Topic1ExplanationsPart1 — pp. 35–39
% Topic 1-9 Questions 1–2, Topic 1-10 Question 1

\needspace{12\baselineskip}
\lesson{Lesson 1-9: Divide Rational Numbers}

\question{Question 1}

\blockhead{Question}
A business loses \$450.75 as a result of a shipping delay. The 3 owners of the
business have to share the loss equally. Write the quotient that represents each
person's share of the loss.

\checked{$-150.25$ (or $-\$150.25$ per owner)}

\blockhead{What the question is asking}
This question asks for each owner's equal share of a \$450.75 loss as a division
answer that keeps the loss sign (negative).

\blockhead{Important words}
\begin{words}
  \item A \term{decimal}uses a decimal point to show parts of a whole.
  \item The \term{decimal point}is the dot that separates whole-number places from parts of a whole.
  \item A \term{loss}is money going out, so we write it as a \term{negative}number.
  \item To share \term{equally}means \term{divide}by the number of people.
  \item A \term{quotient}is the answer to a division problem.
  \item \term{Long division}is a written way to divide one digit at a time.
  \item A \term{remainder}is the amount left after one division step.
  \item To \term{bring down}a digit means move the next digit beside the remainder so division can continue.
  \item In \term{decimal division}\unskip, place the decimal point in the answer directly above the decimal point in the number being divided.
  \item Division sign rule: opposite signs make a negative quotient; same signs make a positive quotient.
\end{words}

\blockhead{Work the problem one tiny step at a time}
\begin{steps}
  \item Write the loss as $-450.75$.
  \item There are 3 owners, so divide: $(-450.75) \div 3$.
  \item Ignore the sign first and divide $450.75 \div 3$.
  \item 3 goes into 4 one time (3), remainder 1.
  \item Bring down 5 to make 15. $3 \times 5 = 15$. Remainder 0.
  \item Bring down the 0 from 450. 3 goes into 0 zero times, so write 0 as the next quotient digit.
  \item The next mark in 450.75 is the decimal point. Put a decimal point in the quotient directly above it.
  \item Bring down 7. 3 goes into 7 two times (6), remainder 1.
  \item Bring down 5 to make 15. $3 \times 5 = 15$. Remainder 0.
  \item So $450.75 \div 3 = 150.25$.
  \item Opposite signs: the quotient is negative. Each share is $-150.25$.
\end{steps}

\finalans{$-150.25$ (or $-\$150.25$ per owner)}

\blockhead{Why this makes sense}
Three equal shares of a loss stay negative. Check: $(-150.25) \times 3 = -450.75$.

\question{Question 2}

\blockhead{Question}
A race is \mixed{15}{1}{2} miles long. The organizers want to put a water station
every \mixed{3}{7}{8} miles along the route, and one at the finish line. How many
total water stations are needed?

\checked{4 water stations}

\blockhead{What the question is asking}
This question asks how many water stations are needed if one is placed every
\mixed{3}{7}{8} miles on a \mixed{15}{1}{2}-mile race, including the finish.

\blockhead{Important words}
\begin{words}
  \item \term{Nonzero}means not equal to zero.
  \item An \term{operation}is a math action. Add joins amounts. Subtract takes away or finds a difference. Multiply makes equal groups. Divide splits into equal groups.
  \item An \term{expression}is a math phrase made of numbers and operation signs.
  \item A \term{whole number}has no fractional part, such as 15.
  \item A \term{fraction}has a top number and a nonzero bottom number and names equal parts of a whole.
  \item The \term{numerator}is the top number of a fraction.
  \item The \term{denominator}is the bottom number of a fraction; it tells how many equal parts make one whole.
  \item A \term{factor}is a number being multiplied.
  \item A \term{mixed number}has a whole number and a fraction, like \mixed{15}{1}{2}.
  \item An \term{improper fraction}has a numerator larger than or equal to the denominator, like $\frac{31}{2}$.
  \item To change a mixed number $a\frac{b}{c}$ to improper: $\dfrac{a \cdot c + b}{c}$.
  \item The \term{reciprocal}of a fraction flips the numerator and the denominator (swap top$\leftrightarrow$bottom). Example: the reciprocal of $\frac{31}{8}$ is $\frac{8}{31}$.
  \item To \term{flip}a fraction means swap numerator and denominator.
  \item To \term{divide fractions}, multiply by the reciprocal: $\dfrac{a}{b} \div \dfrac{c}{d} = \dfrac{a}{b} \times \dfrac{d}{c}$.
  \item To \term{cancel matching factors}means remove the same factor from a numerator and a denominator before multiplying. Example: a 31 on top and a 31 on bottom cancel to 1.
  \item An \term{interval}is one spacing gap between stations.
  \item To \term{simplify}means rewrite a math expression in a shorter form with the same value.
\end{words}

\blockhead{Work the problem one tiny step at a time}
\begin{steps}
  \item Convert race length: $\mixed{15}{1}{2} = \dfrac{15 \times 2 + 1}{2} = \dfrac{31}{2}$.
  \item Convert spacing: $\mixed{3}{7}{8} = \dfrac{3 \times 8 + 7}{8} = \dfrac{31}{8}$.
  \item Divide distance by spacing to count intervals: $\dfrac{31}{2} \div \dfrac{31}{8}$.
  \item Write the reciprocal of the second fraction: flip $\dfrac{31}{8}$ (swap top and bottom) to get $\dfrac{8}{31}$.
  \item Multiply by that reciprocal: $\dfrac{31}{2} \times \dfrac{8}{31}$.
  \item Cancel the matching 31s (same factor top and bottom): $\dfrac{31}{2} \times \dfrac{8}{31}$ becomes $\dfrac{1}{2} \times \dfrac{8}{1}$.
  \item Multiply: $\dfrac{1}{2} \times \dfrac{8}{1} = \dfrac{8}{2}$.
  \item Simplify: $\dfrac{8}{2} = 4$.
  \item That means 4 equal intervals along the race.
  \item Stations are at the end of each interval, including the finish line. So there are 4 stations.
\end{steps}

\finalans{4 water stations}

\blockhead{Why this makes sense}
Four intervals of \mixed{3}{7}{8} fit exactly into \mixed{15}{1}{2}. One station at
the end of each interval gives 4 stations.

\lesson{Lesson 1-10: Solve Problems with Rational Numbers}

\question{Question 1}

\blockhead{Question}
Two people are hiking. The first person climbs to the top of a peak that is 450
feet above sea level. The second person descends into a valley that is 120 feet
below sea level. What is the total vertical distance between the two hikers?

\checked{570 feet}

\blockhead{What the question is asking}
This question asks how many feet separate the two hikers when one is above zero
and the other is below zero.

\blockhead{Important words}
\begin{words}
  \item A \term{negative number}is less than zero. A \term{positive number}is greater than zero.
  \item A \term{number line}is a straight line with numbers in order. Negative numbers are below or left of zero; positive numbers are above or right of zero.
  \item \term{Above sea level}is positive. \term{Below sea level}is negative.
  \item \term{Vertical}means up and down (not side to side).
  \item \term{Vertical distance}is how far apart two heights are on an up-down number line.
  \item \term{Sea level}is 0 feet (the zero mark between above and below).
  \item \term{Absolute value}means distance from zero on a number line. Distance is never negative.
  \item A \term{difference}is the answer to subtraction. To \term{subtract}means take away or find how far apart two numbers are.
  \item A distance can be found with \abs{a - b}: subtract the two positions, then take absolute value. Subtracting a negative adds: $a - (-b) = a + b$.
\end{words}

\blockhead{Work the problem one tiny step at a time}
\begin{steps}
  \item Write the positions: $+450$ (peak) and $-120$ (valley).
  \item Distance $=$ \abs{450 - (-120)}.
  \item Rewrite the subtraction: $450 - (-120) = 450 + 120$.
\end{steps}

\begin{center}
\begin{tikzpicture}[y=0.011cm,x=1cm]
  % vertical number line
  \draw[{Stealth[length=2.4mm]}-{Stealth[length=2.4mm]}] (0,-250) -- (0,500);
  % ticks and labels
  \foreach \y/\lab/\nm in {450/$450$/labyPeak, 225/$225$/labyMid, 0/sea level: $0$/labyZero, -120/$-120$/labyValley, -200/$-200$/labyLow}{%
    \draw (-0.08,\y) -- (0.08,\y);
    \node[right=2pt,font=\footnotesize,inner xsep=2pt] (\nm) at (0.08,\y) {\lab};}
  % marked points
  \fill[red] (0,450) circle (2.2pt);
  \fill[blue] (0,-120) circle (2.2pt);
  \node[left=6pt,font=\footnotesize] at (-0.08,450) {peak};
  \node[left=6pt,font=\footnotesize] at (-0.08,-120) {valley};
  % measurement arrow, clear of the "sea level: 0" label
  \draw[teal,{Stealth[length=2.2mm]}-{Stealth[length=2.2mm]}] (2.35,450) -- (2.35,-120);
  \node[right=3pt,font=\footnotesize,teal] at (2.35,165) {570 ft};
  \draw[teal,dotted] (labyPeak.east) -- (2.35,450);
  \draw[teal,dotted] (labyValley.east) -- (2.35,-120);
\end{tikzpicture}
\end{center}

\begin{steps}[resume]
  \item Add the ones: $0 + 0 = 0$.
  \item Add the tens: $5 + 2 = 7$, so $50 + 20 = 70$.
  \item Add the hundreds: $400 + 100 = 500$.
  \item Total: $500 + 70 = 570$.
  \item Absolute value of 570 is still 570.
  \item Think: 450 feet up from sea level plus 120 feet down from sea level is 570 feet apart.
\end{steps}

\finalans{570 feet}

\blockhead{Why this makes sense}
One hiker is above zero and one is below zero, so you add the two distances from
sea level: $450 + 120 = 570$. Check: \abs{450 - (-120)} $=$ \abs{570} $= 570$.
